\documentclass[11pt]{article}
\usepackage[T1]{fontenc}
\usepackage[utf8]{inputenc}
\usepackage[margin=1in]{geometry}
\usepackage{amsmath,amssymb,amsthm,booktabs,longtable}
\usepackage[colorlinks=true,linkcolor=black,citecolor=blue,urlcolor=blue]{hyperref}
\newtheorem{theorem}{Theorem}[section]
\newtheorem{proposition}[theorem]{Proposition}
\newtheorem{lemma}[theorem]{Lemma}
\newtheorem{corollary}[theorem]{Corollary}
\theoremstyle{definition}
\newtheorem{definition}[theorem]{Definition}
\theoremstyle{remark}
\newtheorem{remark}[theorem]{Remark}
\newcommand{\NN}{\mathbb Z_{\geq0}}
\newcommand{\DD}{D^{2}}
\title{Squares with Square Second Arithmetic Derivative:\\
Infinite Families and an Exact Extension of OEIS A348425}
\author{Carlo Corti\\Independent researcher\\
\href{mailto:carlo.corti@outlook.com}{\texttt{carlo.corti@outlook.com}}}
\date{October 2026}
\begin{document}
\maketitle
\begin{abstract}
OEIS A348425 consists of the nonnegative squares whose second arithmetic
derivative is a square. We consolidate 9,114 consecutive terms, compared
with the 33-term public snapshot acquired on October 22, 2026. The final
campaign contributes 8,917 terms and completes the input-root interval
through $10^{10}$. It follows a 197-term prefix whose entire root interval
has been separately checked. Its arithmetic
kernel uses $D^{2}(r^{2})=rD(r)+2D(r)^{2}+2rD(D(r))$, reducing the irregular
factorization operand to $D(r)$. Complete coverage rests on this exhaustive
producer and its sealed records; separate arithmetic checks verify all
reported positive witnesses and bounded windows. No independent full
re-enumeration of the final campaign interval is claimed. We prove a
necessary and sufficient square criterion for each fixed reduced core and
integer-rate multiplier class, using the historical classification of
$D(t)=\lambda t$. Four families established in earlier project notes are
recovered, and 31 further certified classes are extracted from the existing
data. These constructions give unconditional infinitude without assumptions
on primes represented by quadratic polynomials. An adjacent witness already
present in an acquired prime table sharpens the next-term bound to
$10^{20}<a(9115)\leq100196144067026155201$; it does not determine that term.
\end{abstract}
\noindent\textit{2020 Mathematics Subject Classification:} Primary 11A25;
Secondary 11A41, 11Y55.\\
\textit{Key words and phrases:} arithmetic derivative, integer sequences,
perfect squares, integer-rate multipliers, exact enumeration.

\section{Introduction}\label{sec:intro}
The arithmetic derivative is determined by $D(p)=1$ for every prime $p$,
$D(0)=D(1)=0$, and the product rule
\begin{equation}\label{eq:leibniz}
 D(uv)=uD(v)+vD(u).
\end{equation}
Its factorization formula and elementary bounds are classical
\cite{barbeau,ufn}; a recent exposition is given by Tossavainen et al.\
\cite{survey}. Throughout this paper, $D^{2}(n)=D(D(n))$ denotes composition,
whereas $D(n)^{2}$ denotes an ordinary square. The function $D$ is not additive.

Marius A. Burtea introduced OEIS A348425 in October 2021 \cite{oeis}:
\begin{equation}\label{eq:set}
 \mathcal A=\{r^{2}:r\in\NN,\ D^{2}(r^{2})\text{ is a square in }\NN\}.
\end{equation}
We write $a(1),a(2),\ldots$ for its increasing enumeration. Thus
$a(1)=0$, $a(2)=1$, $a(3)=4$, and $a(4)=49$ (offset $1$); for example,
$D(49)=14$ and $D(14)=9$. The input root $r$, the sequence index $i$, and
the second-derivative root $R=\sqrt{D^{2}(r^{2})}$ are different quantities.

The snapshot acquired on October 22, 2026 contains 33 indexed terms.
Its last value is $31618129=5623^{2}$. It also contains root-enumeration programs by
Amiram Eldar and Michel Marcus and three prime-power constructions
\cite{oeis}. We retain that precise baseline when reporting an extension;
the final campaign alone did not produce all 9,081 terms beyond it.

There are two complementary questions. Exact indexing requires checking
every intervening input root. A structural family supplies infinitely many
members, but cannot establish their sequence indices without that coverage.
The first question is addressed by the completed enumeration and its
validation. For the second, a product-rule criterion transfers the historical
integer-rate multipliers of Kovi\v{c} \cite[Proposition 2]{kovic} to square
second derivatives. The corresponding reduced-core decomposition is already
present in Ufnarovski and \AA hlander \cite[Section 8]{ufn}.

The contributions are an exact local extension to 9,114 terms; proofs and
disjointness statements for scale families; a reproducible finite analysis
of prime powers and reduced cores; and a tighter, unindexed witness bound.
The four initially supplied scale families predate the final campaign:
their mechanism was developed in project notes on October 5, 2026, with
the fourth seed retained in subsequent notes. The
31 additional classes below are identified by analyzing the consolidated
data. We make no claim of absolute bibliographic priority for these
specializations. In particular, the book of Balzarotti and Lava
\cite{book} has not been consulted in full.

\section{Exact identities and scale families}\label{sec:math}
For $n\geq2$ with prime factorization $n=\prod_p p^{e_p}$, the defining
rules give
\begin{equation}\label{eq:factor}
 D(n)=\sum_{p\mid n}e_p\frac{n}{p}.
\end{equation}
We use the empty product convention for $1$ and treat $0$ separately.

\begin{proposition}[Root reduction]\label{prop:root}
For $r\in\NN$, put $d=D(r)$ and $e=D(d)$. Then
\begin{equation}\label{eq:root}
 D(r^{2})=2rd,\qquad \DD(r^{2})=rd+2d^{2}+2re.
\end{equation}
\end{proposition}
\begin{proof}
The first identity follows from \eqref{eq:leibniz}. For the second,
differentiate the product $2(rd)$, using $D(2)=1$:
\[
 D(2rd)=rd+2D(rd)=rd+2\bigl(d^{2}+rD(d)\bigr).
\]
This uses the product rule twice; it does not differentiate a sum
term by term. Both identities hold also for $r=0,1$.
\end{proof}

For prime $r$, the result is $\DD(r^{2})=r+2$. Prime input roots are
therefore essential, and may not be discarded. More generally, the
identity replaces factorization of $2rD(r)$ by factorization of $D(r)$.

\begin{lemma}[Integer-rate multipliers]\label{lem:rate}
Let $\lambda\in\NN$. The positive integer solutions of $D(t)=\lambda t$
are exactly
\begin{equation}\label{eq:mult}
 t=\prod_p p^{p c_p},\qquad c_p\in\NN,\qquad \sum_p c_p=\lambda,
\end{equation}
where the support is finite. In particular, $\lambda=0$ forces $t=1$.
\end{lemma}
\begin{proof}
This is the integer specialization of the classical logarithmic
decomposition \cite{ufn} and Kovi\v{c}'s Proposition 2 \cite{kovic}.
The case $t=1$ has $D(t)=0$ and $\lambda=0$, and satisfies the
empty-product convention. For $t>1$, write $t=\prod_{p\mid t}p^{e_p}$. Then
$\sum_{p\mid t}e_p/p=\lambda$. Multiply by $P=\prod_{p\mid t}p$ and reduce
modulo a fixed prime $p\mid t$. All terms except $e_pP/p$ vanish, so
$p\mid e_p$. Setting $e_p=pc_p$ proves necessity; substituting in
\eqref{eq:factor} proves sufficiency. For $\lambda=0$, all $c_p=0$, so $t=1$.
\end{proof}

Equivalently, $t$ is a product of $\lambda$ factors $p^{p}$, with primes
allowed to repeat. There is no coprimality restriction between these
primes and any separately chosen seed.

\begin{theorem}[Square criterion in a scale class]\label{thm:scale}
Let $m,t\geq1$ and $\lambda\in\NN$ satisfy $D(t)=\lambda t$. Define
\begin{equation}\label{eq:coefficient}
 b=D(m)+\lambda m,\qquad
 C(m,\lambda)=2D(mb)+(4\lambda+1)mb.
\end{equation}
Then
\begin{equation}\label{eq:scale}
 \DD((mt)^{2})=t^{2}C(m,\lambda).
\end{equation}
Consequently $(mt)^{2}\in\mathcal A$ if and only if $C(m,\lambda)$
is a square. If $C(m,\lambda)=S^{2}$, its second-derivative root is $St$.
\end{theorem}
\begin{proof}
The product rule gives $D(mt)=tb$, hence $D((mt)^{2})=2t^{2}mb$.
Moreover, $D(t^{2})=2\lambda t^{2}$ and
$D(2t^{2})=(4\lambda+1)t^{2}$. Differentiating the product
$(2t^{2})(mb)$ proves \eqref{eq:scale}. The coefficient is a nonnegative
integer. Multiplication by $t^{2}$ preserves every prime-exponent parity,
so it is a square exactly when the product is. The only zero case is
$m=1$, $\lambda=0$, hence $t=1$, giving the member $1$.
\end{proof}

\begin{definition}\label{def:core}
For $r\geq1$, write each prime exponent uniquely as
$e_p=pc_p+b_p$, with $0\leq b_p<p$, and set
\begin{equation}\label{eq:canonical}
 m=\prod_p p^{b_p},\qquad t=\prod_p p^{pc_p},\qquad
 \lambda=\sum_p c_p.
\end{equation}
We call $m$ the reduced core and $\lambda$ the canonical rate of $r$.
A positive integer $m$ is \emph{reduced} when $\nu_p(m)<p$ for every
prime $p$, where $\nu_p$ is the prime exponent.
\end{definition}

This is the small/large-parts decomposition of \cite[Section 8]{ufn},
expressed in input-root
notation. It is unique, $r=mt$, and Lemma~\ref{lem:rate} applies to $t$.
Thus Theorem~\ref{thm:scale} is an exact criterion for each core and rate,
but does not solve all possible conditions $C(m,\lambda)=S^{2}$.

\begin{corollary}\label{cor:infinite}
If $m$ is reduced, $\lambda\geq1$, and $C(m,\lambda)$ is a square, then
\[
 \left\{\left(m\prod_{j=1}^{\lambda}p_j^{p_j}\right)^{2}:
       p_1,\ldots,p_\lambda\text{ prime}\right\}\subseteq\mathcal A
\]
is infinite. Families from distinct reduced-core/rate pairs are disjoint.
\end{corollary}
\begin{proof}
Membership follows from the preceding theorem and lemma. Fix all but one
prime parameter, and let the remaining prime increase; the roots are
unbounded. For disjointness, multiplication adds multiples of $p$ to
the exponent of each prime $p$, even when $p$ divides $m$. Since
$0\leq\nu_p(m)<p$, adding these multiples produces no carry from
the reduced exponent into the quotient. The residues
and the total quotient sum in \eqref{eq:canonical} recover $m$ and
$\lambda$ uniquely. A root cannot belong to two distinct pairs.
\end{proof}

\subsection{Four explicit scale families}\label{sec:four}
Table~\ref{tab:four} gives exact certificates for the four seeds from
the earlier notes.
Every seed is reduced: they are $47$, $5^{3}$, $5^{3}7$, and $23^{2}$.
They therefore yield four disjoint infinite families:
\begin{align*}
 (47p^{p})^{2},&\quad R=148p^{p},&
 (125p^{p})^{2},&\quad R=500p^{p},\\
 (875p^{p})^{2},&\quad R=3125p^{p},&
 (529p^{p}q^{q})^{2},&\quad R=2852p^{p}q^{q}.
\end{align*}
Here $p,q$ are arbitrary primes, including $p=q$, and may divide the
seed. No conjecture about polynomial primes is involved.

\begin{table}[htbp]
\centering
\caption{Exact certificates for the inherited families. The notation is
that of \eqref{eq:coefficient}; $C=S^{2}$.}\label{tab:four}
\begin{tabular}{rrrrrr}
\toprule
$m$&$\lambda$&$b$&$D(mb)$&$C$&$S$\\
\midrule
47&1&48&5312&21904&148\\
125&1&200&62500&250000&500\\
875&1&1525&1546875&9765625&3125\\
529&2&1104&1438880&8133904&2852\\
\bottomrule
\end{tabular}
\end{table}

Formulations using the square seed $m^{2}$ describe the same families:
$D(t^{2})=2\lambda t^{2}$, so the corresponding multiplier rate for
square inputs is $2\lambda$, rather than $\lambda$.

\subsection{Prime powers and the parameter convention}\label{sec:powers}
\begin{proposition}\label{prop:power}
For a prime $p$ and $k\geq1$,
\begin{equation}\label{eq:power}
 \DD(p^{2k})=p^{2k-2}\bigl[pD(2k)+2k(2k-1)\bigr].
\end{equation}
Thus $p^{2k}\in\mathcal A$ exactly when the bracket is a square.
\end{proposition}
\begin{proof}
Differentiate $D(p^{2k})=(2k)p^{2k-1}$ as a product. The result is
$D(2k)p^{2k-1}+(2k)(2k-1)p^{2k-2}$. The remaining factor
$p^{2k-2}$ is a square. This calculation does not assume $p\nmid2k$.
\end{proof}

For $k=1,2,3$, the respective brackets are
\begin{equation}\label{eq:three}
 p+2,\qquad 4(p+3),\qquad 5(p+6).
\end{equation}
These recover the three constructions in A348425 \cite{oeis}:
$p=s^{2}-2$, $p=s^{2}-3$, and $p=5s^{2}-6$, respectively, subject
to primality of $p$. They do not prove that the required prime sets
are infinite. In particular the third condition follows from
$5(p+6)=u^{2}$: $5\mid u$, so $u=5s$ and $p=5s^{2}-6$.

The parameter conventions of the related entries differ
\cite{near2,near3,near5}: A028871 lists the primes $p=s^{2}-2$;
A028873 lists the parameters $s$ for which $s^{2}-3$ is prime;
A201787 lists the primes $p=5s^{2}-6$, not their parameters.
The third comment in the checked A348425 snapshot uses an A201787 term
as if it were $s$. For its second term, $599$, the literal assertion
fails because
\[
 5\cdot599^{2}-6=1793999=19\cdot94421.
\]
The corrected mathematical statement is: if the prime $p\in\mathrm{A201787}$
has $p=5s^{2}-6$, then $\DD(p^{6})=(5p^{2}s)^{2}$. The live text
checked on October 22, 2026 retains that parameter error.

\section{Enumeration, arithmetic bounds, and completion}\label{sec:algorithm}
Since $r\mapsto r^{2}$ is strictly increasing on $\NN$, exhaustive
root enumeration preserves the order in \eqref{eq:set}. The producer
first computes $D(r)$ in blocks, then applies Proposition~\ref{prop:root}
and an exact integer-square test to every root. No family restriction,
primality restriction on $r$, or empirical exclusion is used.

\subsection{Sieve and exact arithmetic}
For a block $[L,H]$, initialize a residual equal to each root and a zero
derivative accumulator. For each prime $p\leq\sqrt H$, strip $p^{e}$
from the residual and add $e(r/p)$ using the original $r$. Any final
residual greater than $1$ is prime, since a composite residual would
have a prime factor at most $\sqrt H$ already processed. The residual
prime $q$ has exponent one: $q>\sqrt H$ implies $q^{2}>H$. Its contribution
is $r$ divided by that residual. This proves that the sieve gives
$d=D(r)$ exactly. The cases $r=0,1$ are handled explicitly.

The irregular operand $d$ is factored using small-prime division,
deterministic word primality based on the Miller--Rabin test
\cite{miller,rabin}, and a fixed-seed rho accelerator of
Pollard type \cite{pollard}. This uses Floyd cycle detection,
constants $c=1,\ldots,32$, and at most 100,000 steps per constant.
The primality bases are
\[
 2,\ 325,\ 9375,\ 28178,\ 450775,\ 9780504,\ 1795265022,
\]
whose coverage for integers below $2^{64}$ is the Sinclair record
listed in \cite{mr}. The Feitsma--Galway tables \cite{feitsma} supply
exhaustive base-2 pseudoprime data for this word range; these large
tables were not independently regenerated here. Bases are reduced
modulo the candidate and a zero residue is skipped. A finite check
of all composite divisors of these bases confirms rejection by the
remaining tests, including $14089=73\cdot193$, which survives the
initial small-prime screen. Every factor leaf is checked for primality, and
the product of the prime powers must reproduce the operand exactly.
Exhausting the bounded rho attempts invokes complete odd trial division;
it never turns a factorization failure into a negative result. This
fallback ensures exactness, without a uniform practical running-time
claim for adversarial operands.

The sieve supplies $d$; only $d$ needs irregular factorization for
each root. The producer does not factor $e=D(d)$. After obtaining $e$,
it computes the two expressions
in \eqref{eq:root}. The residue condition
$D^{2}(r^{2})\bmod16\in\{0,1,4,9\}$ is necessary only. A floating
square-root estimate is corrected with integer products until
$q^{2}\leq y<(q+1)^{2}$, where $y=D^{2}(r^{2})$; acceptance requires
$q^{2}=y$. The decision is therefore integral, rather than a
floating-point test.

\begin{lemma}[Width bounds]\label{lem:width}
For $U\geq2$, let $F(U)=\lfloor U\lfloor\log_2U\rfloor/2\rfloor$,
with $F(0)=F(1)=0$. If $0\leq v\leq U$, then $D(v)\leq F(U)$.
For $r\leq B$, set $A=F(B)$ and $E=F(A)$. Then
\begin{equation}\label{eq:width}
 r^{2}\leq B^{2},\quad d\leq A,\quad e\leq E,\quad
 D(r^{2})\leq2BA,\quad
 \DD(r^{2})\leq BA+2A^{2}+2BE.
\end{equation}
\end{lemma}
\begin{proof}
For $v\geq2$, \eqref{eq:factor} gives $D(v)\leq v\Omega(v)/2$, where
$\Omega(v)=\sum_pe_p$. Since $2^{\Omega(v)}\leq v$,
$\Omega(v)\leq\lfloor\log_2v\rfloor$. The integer bound follows,
because $v\lfloor\log_2v\rfloor$ is nondecreasing for integer $v\geq2$;
the cases $v=0,1$ follow separately. Applying it first to $r$ and then
to $d$, and using \eqref{eq:root}, proves \eqref{eq:width}.
This is the integer form of the classical upper bound \cite{barbeau}.
\end{proof}

The supported source capacity is $B=10^{12}$. Here
\[
 A=19500000000000,\qquad E=429000000000000,
\]
\[
 B^{2}=10^{24},\quad 2BA=3.9\cdot10^{25},\quad
 BA+2A^{2}+2BE=1.638\cdot10^{27}<2^{91}.
\]
Thus roots and factorization operands fit unsigned 64-bit words, while
coordinates and their products use unsigned 128-bit arithmetic, with
promotion before multiplication. The second-derivative root is at most
$40472212689696$. These are conservative capacity bounds, not
measurements or coverage through $10^{12}$.

\subsection{Ordered completion and the executed identity}\label{sec:provenance}
The completed campaign used Linux Mint 22.3 on x86-64 LP64,
GCC 13.3.0, OpenMP with 16 workers, and blocks of 65,536 roots.
The processor was an AMD Ryzen 9 7940HS (8 cores, 16 hardware threads),
with 64 GB DDR5 installed RAM. The runtime libraries were GNU
libgomp (package 14.2.0-4ubuntu2\textasciitilde24.04.1) and OpenSSL
3.0.13; the Python checks used Python 3.12.3 and SymPy 1.12.
Appendix~\ref{app:records} gives the exact build options and record grammar. Within a block each worker writes its own result
flag. Positive rows are replayed in increasing root order after the
join. A block advances the committed frontier only when all its roots
finish; an interrupted or failed block cannot certify negatives.
The two sealed segments in Table~\ref{tab:segments} are contiguous.

\begin{table}[htbp]
\centering
\caption{Completed final-campaign segments, with inclusive input roots.
The hit counts sum to 8,917.}\label{tab:segments}
\begin{tabular}{lrrr}
\toprule
Segment&First root&Last root&Hits\\
\midrule
0&1000001&5454971456&6792\\
1&5454971457&10000000000&2125\\
\bottomrule
\end{tabular}
\end{table}

The author executed the campaign with the historical source
\texttt{a348425\_v2.c}. The public-facing copy is
\texttt{a348425.c}. Its sole C change is the neighboring-source filename
used for identity checking; the Makefile changes that filename and
verification paths. The arithmetic, capacity, parallelism, record grammar,
and build options are unchanged. The original executed source, Makefile,
binary, 197-term baseline, and saved run records are preserved separately,
with an exact transformation record. The normalized build is not credited
with producing the historical results. These provenance materials belong
to the locally prepared project record \cite{project}; they do not imply
that a release has been published.

The sum of saved internal elapsed counters is approximately 48 minutes
46 seconds. It excludes part of checkpoint I/O and other startup or
closure costs and is not the complete command wall time. A campaign
peak-RSS measurement is unavailable; short calibration measurements
are not substituted for it. Rho-fallback counts were printed during
execution but are not retained in the saved campaign records, so
no observed fallback frequency is reported. No unmeasured performance
comparison is made.

\section{Indexed extension and its validation}\label{sec:results}
\begin{theorem}[Computational result]\label{thm:computed}
The separately checked prefix and the completed exhaustive producer
give 9,114 consecutive terms of $\mathcal A$ through $r=10^{10}$.
Coverage of the final interval $1000001\leq r\leq10^{10}$ rests on
this producer and its preserved completion records, rather than an
independent re-enumeration of that full interval.
In particular,
\begin{align*}
 a(36)&=54582544=7388^{2},\\
 a(198)&=1012050084049=1006007^{2},\\
 a(9114)&=99948010099131627889=9997400167^{2}.
\end{align*}
No further member has input root at most $10^{10}$. The first unresolved
index is 9115.
\end{theorem}
\begin{proof}[Computational evidence and coverage argument]
The chronology is given in Table~\ref{tab:stages}. The original extension
admitted the 33-term public baseline. During manuscript revision on
October 23, 2026, complete trial division applied directly to $r^{2}$
and $D(r^{2})$ for every $0\leq r\leq6500$ found exactly the 35
listed members, including exactly those 33 through $r=5623$. A separately
supplied C implementation agreed on this full interval; it factors
$r$ and $D(r)$ and evaluates $D(2rD(r))$ from the merged prime multiset.
This closes the earlier admitted-prefix dependency, without attributing
the later check to the original campaign. Before that revision, an exact replay of
877 roots, $5624\leq r\leq6500$, found the two additional roots
$5625$ and $5927$. Direct trial division and a separate SymPy
computation agreed on this whole interval. A further complete finite
check of 993,500 roots, $6501\leq r\leq10^{6}$, gave 162 more members,
agreeing with a separate completed producer run on that finite
interval and yielding a 197-term prefix. This check
used a separately supplied Python implementation: it factors $r$
and $D(r)$ by trial division and forms the prime multiset of
$2rD(r)$, rather than evaluating the second formula in \eqref{eq:root}.
The delivered implementation and its local execution record are retained.

The final producer covers exactly $10^{10}-10^{6}=9999000000$ further
roots in the two segments of Table~\ref{tab:segments}. Its committed
cursors reach the inclusive upper endpoints, the segments join without
gaps or overlap, and the sealed hit lists agree with their ordered
aggregate. They contain $6792+2125=8917$ positives. The preserved prefix
and these ordered hits give $197+8917=9114$ indexed rows. The root
reduction, exact tests, and complete-block commitment establish how the
producer accounts for negatives. In particular, the last positive root
is not used as a proxy for the completed frontier.
\end{proof}

\begin{table}[htbp]
\centering
\caption{Distinct stages of the local extension, relative to the acquired
public baseline. Its completeness was separately checked on October 23,
2026; this later verification adds no terms. The extension beyond the
snapshot is $9114-33=9081$ terms.}\label{tab:stages}
\begin{tabular}{lrrr}
\toprule
Stage&Root frontier&Added terms&Total terms\\
\midrule
Public snapshot&5623&---&33\\
Exact replay&6500&2&35\\
Complete finite continuation&1000000&162&197\\
Final campaign&10000000000&8917&9114\\
\bottomrule
\end{tabular}
\end{table}

\subsection{What the separate checks establish}\label{sec:validation}
All 8,917 final-campaign positives were recomputed in Python with
SymPy 1.12 \cite{sympy}, factoring $n=r^{2}$ and $D(n)$ directly.
Factor products and prime-leaf domains were checked. A documentary
table with columns
\[
 (i,\ r,\ n=a(i),\ D(n),\ \DD(n),\ R)
\]
was then checked for every one of its 9,114 rows: 18,228 derivative
identities and 18,228 square identities, together with equality to the
indexed b-file. This table records witnesses for the contiguous terms;
it is not a sparse extension across a missing index.

Four complete windows were also checked with separate arithmetic:
\begin{gather*}
 [1000001,1001419],\quad[100000001,100001000],\\
 [1000000001,1000001000],\quad[9999999001,10000000000].
\end{gather*}
They contain $1419+1000+1000+1000=4419$ distinct roots, all negative,
in agreement with the run. Positive-path agreement is supplied by
the all-positive check. During manuscript revision, the separately
supplied C checker also reproduced complete 129-root windows
$[1005943,1006071]$ and $[9997400103,9997400231]$, containing
respectively the first and last campaign hits. A third 129-root window
$[5454971392,5454971520]$ straddles the segment join and is negative.
These bounded checks supplement, and do not extend, campaign coverage.

These checks have complementary limitations. Recomputing a reported
positive does not detect an omitted positive. Seals, identity checks,
aggregate equality and completed cursors establish integrity and the
producer's recorded completion, rather than a second independent
enumeration of all 9,999,000,000 final-campaign inputs. The 33-term
baseline dependency was closed by the complete check described above.
The separately
implemented finite checker shares the mathematical definition, and
implementations supplied by one LLM family and reviews from that
family have correlated provenance. Execution claims in external code reviews that
lack delivered raw harnesses or logs are not counted as local reruns.

Source review before execution concerns software correctness and
does not by itself verify the structural theorems. During manuscript
preparation, the current data were
again checked by the finite-analysis companion, and the archived
original run was verified without advancing it. The provenance
checker also regenerated all 9,114 indexed pairs from the original
197-term baseline and the sealed final hit list.

\subsection{The first unresolved term}\label{sec:bound}
The families already show that $\mathcal A$ is infinite, so the next
term exists. The completed negative tail gives $a(9115)>10^{20}$.
The earlier family witness
\[
 r=529\cdot7^{7}\cdot3^{3}=11762664669
\]
has square $138360280115340879561$ and second-derivative root
$63416105172$, giving a valid unindexed upper bound.

Inspection of the acquired A028871 table \cite{near2} during manuscript
preparation on October 22, 2026 identified a tighter witness, also
present in the earlier restricted-family list:
\[
 p=10009802399=100049^{2}-2.
\]
Complete trial division through $\lfloor\sqrt p\rfloor=100048$
confirms primality. Proposition~\ref{prop:root} then gives
\[
 D(p^{2})=20019604798,\qquad \DD(p^{2})=10009802401=100049^{2}.
\]
Consequently
\begin{equation}\label{eq:bound}
 10^{20}<a(9115)\leq100196144067026155201.
\end{equation}
The provider index of $p$ in A028871 is 8889; it is not an A348425
index. No arbitrary-root interval beyond $10^{10}$ was enumerated
to obtain this improvement. In particular, the witness is not asserted
to equal $a(9115)$.

\section{Finite structure of the consolidated data}\label{sec:finite}
The analysis in this section examines only the existing 9,114 positive
rows. Its companion program checks direct composition, factors each
input root, constructs \eqref{eq:canonical}, and verifies the exact
coefficient and root in Theorem~\ref{thm:scale}. It searches neither
additional roots nor a larger seed domain.

\subsection{Prime powers and a disjoint classification}
Apart from $r=0,1$, there are 8,888 prime roots, 52 higher prime-power
roots, and 172 roots with at least two distinct prime factors.
The prime-power exponents in $r=p^{k}$ are given in
Table~\ref{tab:exponents}.
\begin{table}[htbp]
\centering
\caption{Prime-power exponents among input roots $2\leq r\leq10^{10}$
of members of $\mathcal A$.}
\label{tab:exponents}
\begin{tabular}{rrrrrr}
\toprule
$k$&1&2&3&8&20\\
\midrule
Count&8888&45&4&2&1\\
\bottomrule
\end{tabular}
\end{table}
The exceptional higher exponents are $5^{8}$, $17^{8}$ and $3^{20}$.
They satisfy Proposition~\ref{prop:power}; $5^{8}$ belongs to the
inherited core $125$, rate $1$ family, and $3^{20}$ to the core $9$,
rate $6$ family below. Thus the three small-exponent constructions
do not exhaust the prime-power part.

As a separate provider check, transforming all acquired terms of
A028871, A028873 and A201787 within the completed root bound gives
respectively 8,888, 45 and 4 members, all present in the consolidated
table. These are restricted-domain checks, not coverage for mixed roots.
The following disjoint partition classifies membership in the
displayed constructions:
\begin{equation}\label{eq:partition}
 9114=2+8937+19+111+45.
\end{equation}
The summands count, in order, $r=0,1$; the $k=1,2,3$ prime-power
constructions; the four inherited scale families; the other observed
positive-rate classes; and the remaining rate-zero members. Here
$8937=8888+45+4$. The inherited classes contain respectively
4, 4, 4 and 7 terms. The classification assigns $r=0,1$ first, then the $k=1,2,3$
prime-power constructions, then the positive-rate classes, and finally
the remaining roots. No summand in \eqref{eq:partition} counts a
term twice. The last 45 terms are not unexplained by the defining
predicate: they are outside this particular collection of families.

\subsection{Canonical rates and additional certified families}
Among the 9,112 roots at least $2$, the rate distribution is given in
Table~\ref{tab:rates}.
\begin{table}[htbp]
\centering
\caption{Canonical rates among the 9,112 input roots at least $2$.}
\label{tab:rates}
\begin{tabular}{rrrrrrrrr}
\toprule
$\lambda$&0&1&2&3&4&5&6&7\\
\midrule
Count&8982&28&29&17&3&27&23&3\\
\bottomrule
\end{tabular}
\end{table}
There are exactly 35 distinct observed core/rate pairs with
$\lambda>0$, accounting for 130 terms. All their reduced cores
are distinct. Table~\ref{tab:cores} records $S=\sqrt{C(m,\lambda)}$
and the finite number $h$ of terms in the completed domain.
Each row gives a universal infinite family by
Corollary~\ref{cor:infinite}, not just its $h$ observed members.
The four marked rows are the inherited families; the other 31 are
additional certified specializations extracted during this finite
analysis, without a bibliographic-first claim.

For example, $m=9$, $\lambda=6$ gives $b=60$ and $D(mb)=D(540)=1188$,
so
\[
 C(9,6)=2\cdot1188+25\cdot540=15876=126^{2}.
\]
Hence
\[
 \left(9\prod_{j=1}^{6}p_j^{p_j}\right)^{2}\in\mathcal A,
 \qquad R=126\prod_{j=1}^{6}p_j^{p_j}
\]
for arbitrary primes, including repetitions and $p_j=3$.
As another example, $m=1847$, $\lambda=1$ has $S=5584$;
its $p=2$ member has root $7388$ and is exactly the newly ranked
$a(36)$. The structural consequence is derived from the data and
the theorem, while the member's rank comes from exhaustive coverage.

\begin{longtable}{rrrr}
\caption{All observed positive-rate reduced-core classes. Each row
satisfies $C(m,\lambda)=S^{2}$. The count $h$ is restricted to
$r\leq10^{10}$; a star marks one of the four inherited families.}
\label{tab:cores}\\
\toprule
$m$&$\lambda$&$S$&$h$\\
\midrule
\endfirsthead
\toprule
$m$&$\lambda$&$S$&$h$\\
\midrule
\endhead
\bottomrule
\endfoot
9&6&126&12\\
34&5&396&11\\
$47^{*}$&1&148&4\\
$125^{*}$&1&500&4\\
$529^{*}$&2&2852&7\\
847&3&7216&7\\
$875^{*}$&1&3125&4\\
1369&2&6586&6\\
1847&1&5584&4\\
2209&5&23171&6\\
6103&6&76862&4\\
6929&7&102544&3\\
8750&5&111750&5\\
28561&1&81289&3\\
33746&6&445510&3\\
41503&2&229075&5\\
230911&6&3076116&2\\
296009&6&3931172&2\\
356377&3&2507085&4\\
371293&5&4227028&2\\
1067089&2&5059634&3\\
1159171&5&13891280&2\\
1830125&1&8984250&3\\
2724449&3&19985875&3\\
3457233&4&33665463&2\\
5182723&1&18966532&2\\
5185474&2&32708088&3\\
7140250&3&63163750&2\\
9020273&5&92443535&1\\
9759362&2&57933234&3\\
25574423&2&146048812&2\\
35701250&4&430062750&1\\
47316959&1&175360104&2\\
62344607&1&252382104&2\\
105971074&3&796214804&1\\
\end{longtable}

The table classifies the positive-rate classes represented in this
finite dataset. It does not classify all seeds or rates. In particular,
the absence of a core from the table does not rule out a family whose
smallest member is outside the completed range. The large prime-root
count and the rate histogram imply neither an asymptotic density nor
an exhaustive list of constructions.

\section{Relations and remaining questions}\label{sec:relations}
The defining relation uses the square inputs of A000290, the derivative
A003415 and its composition A068346 \cite{squares,derivative,second}. The paired tables
A056798 and A192084 have a different meaning: at their common index
$j$, the first gives an even prime power (including $1$) and the
second gives its second derivative \cite{evenpower,powersecond}. Their acquired 10,000 aligned
rows contain 72 square-valued second derivatives, with input table
ending at $10708731289$. Those indices are provider indices; they are
not positions in A348425. Moreover, prime-power tables cannot account
for mixed roots.

A266890 requires the first arithmetic derivative of a square to be
square \cite{firstsquare}. Checking $D(n)$ in every consolidated row,
using the coverage of Theorem~\ref{thm:computed}, gives
the finite intersection
\[
 \mathcal A\cap\mathrm{A266890}\cap[0,10^{20}]=\{0,1,4\}.
\]
This does not assert a global intersection theorem. For example,
$49\in\mathcal A$ but $D(49)=14$ is not square, whereas $256$ has
$D(256)=1024$ square and $D^{2}(256)=5120$ not square.

There are also two unconditional exclusions. For $r\geq2$,
$D(r)\geq1$, and \eqref{eq:root} gives $D^{2}(r^{2})\geq r\geq2$.
If such a value is square, it is at least $4$, hence neither prime
nor squarefree. For $r=0,1$, the value is $0$, also neither prime
nor squarefree. Thus A348425 is disjoint from A192192, whose second
derivative is prime, and A328244, whose second derivative is squarefree
\cite{prime2,sfree2}.

The natural structural problem left by Theorem~\ref{thm:scale} is to
describe reduced cores $m$ and rates $\lambda$ for which
$C(m,\lambda)$ is square, beyond finite tables of certificates.
The prime-power criterion supplies another restricted classification
problem as $k$ varies. Neither problem is resolved by the observed
counts, and no conjectural density model is used here. The computational
question immediately after this work remains the first unresolved
index 9115: establishing it requires arbitrary-root coverage beyond
$10^{10}$, not merely testing another supplied positive witness.

\section{Conclusion}
The consolidated extension contains 9,114 consecutive terms through
input root $10^{10}$, with last term $99948010099131627889$.
The 33-term public snapshot, the two finite continuation stages,
and the 8,917-term final campaign have distinct evidential roles.
A later complete check closes the original public-prefix dependency;
negative coverage of the final interval rests on the single exhaustive
producer and its completion records.
The next term satisfies \eqref{eq:bound}, while its exact value remains
undetermined. The scale criterion turns canonical integer-rate
decomposition into provable infinite families, including 35 disjoint
classes represented in the data. It explains a structured part of the
mixed-root behavior without converting a finite census into a global
classification.

\section*{Data availability}
The indexed terms, documentary witness table, normalized C source and
build instructions, original execution materials, provenance checks,
and Python validators currently exist in a local preparation. The
manuscript and finite-analysis companion are separate local review
materials. No public release of this project is established at the
date of this draft. These materials can be supplied privately for
review; their intended repository and archive are identified in
\cite{project}. Those locators are prospective, and the archive DOI
is an intentional placeholder pending the author's verified publication
locator. Original source reviews remain internal materials. Complete
independent re-enumeration of the final search interval is not supplied.

\section*{AI use disclosure}
Large language models were used as research and drafting assistants during
the A348425 workflow, including OpenAI Codex and external LLM-based
review systems from the ChatGPT, Claude, DeepSeek, Gemini, Grok, Qwen,
MiniMax, Meta AI, GLM and Mistral families. These are the product or family
names present in the project records; exact model-version metadata were
not consistently available or authenticated. The retained history spans
October 5--23, 2026; the first manuscript review was integrated on
October 23. These channels are not claimed to be independent providers.
\paragraph{Mathematical analysis.}
OpenAI Codex assisted with literature research, product-rule derivations,
canonical-core analysis, and the distinction between indexed terms and
unindexed witnesses. The four inherited families were already present
in project notes before the final campaign. The manuscript's additional
class certificates and tightened bound were checked by exact arithmetic
on retained data and primary sources, not by accepting conversational
assertions as evidence.
\paragraph{Algorithm and code.}
OpenAI Codex assisted with drafting the C implementation, validation
planning, and integration of supplied code critiques. Claude-supplied
Python and C implementations were locally rerun for the finite
continuation and bounded revision checks, respectively; they do not
constitute independent full-campaign reproduction. Such assistance and code review are fallible and do not
replace the proofs, executable checks, or the stated coverage boundaries.
The author ran the official campaign. No personal verification of
every proof or manual recalculation of every output row is asserted.
\paragraph{Manuscript.}
OpenAI Codex assisted with drafting, source reconciliation, and internal
checking of this text. The executable validators, result tables,
original run records and primary sources form the verification basis.
Carlo Corti retains full scientific responsibility for the mathematical
statements, implementation, numerical results, validation boundaries,
and conclusions.

\appendix
\section{Build and completion-record specification}\label{app:records}
The executed Makefile uses \texttt{gcc}, with \texttt{-O3 -std=c17},
\texttt{-Wall}, \texttt{-Wextra}, \texttt{-Wpedantic},
\texttt{-Wshadow}, \texttt{-Wconversion},
\texttt{-Wdouble-promotion}, \texttt{-Wformat=2}, \texttt{-Werror},
and \texttt{-fopenmp}; linking uses
\texttt{-fopenmp -lcrypto -lm}. Two preprocessor string definitions bind the source and Makefile identities:
\[\texttt{A348425\_SOURCE\_SHA},\qquad \texttt{A348425\_MAKE\_SHA}.\] The original build file, rather than a
reconstruction of its shell quoting, is retained with the execution materials.

The configuration starts with \texttt{A348425-CFG-1} and records
\texttt{start}, \texttt{end}, \texttt{block}, \texttt{group},
\texttt{threads}, and the source, Makefile and executable bindings.
The executed values are $1000001$, $10^{10}$, $65536$, $2$ and $16$;
\texttt{group} is the number $S$ of segments. More precisely, let
$u,v$ be the configured inclusive root endpoints and $\ell$ the block size.
There are $K=\lfloor(v-u)/\ell\rfloor+1$ blocks. Define
\[
 q_j=j+\left\lfloor\frac{(K-S)j(6S+1-j)}{S(5S+1)}\right\rfloor,
 \qquad 0\leq j\leq S.
\]
Segment $j$, numbered from $0$, covers
$[u+q_j\ell,\,\min(v,u+q_{j+1}\ell-1)]$ for $0\leq j<S$.
This reproduces the boundaries in Table~\ref{tab:segments}; $q_0=0$
and $q_S=K$, and the reserved block in each positive increment
$q_{j+1}-q_j$ ensures nonempty contiguous segments.

A segment starts with \texttt{A348425-SEG-1}, followed by configuration
binding, segment identifier, inclusive \texttt{low} and \texttt{high},
\texttt{next}, internal \texttt{elapsed\_ns}, and \texttt{hits}.
Its positive records have five tab-separated unsigned decimal columns
$(r,r^{2},D(r^{2}),D^{2}(r^{2}),R)$ in strictly increasing $r$.
The trailing \texttt{state\_sha} binds every preceding byte of that segment;
a separate seal binds the full segment, including this trailing line.
Segment filenames are \texttt{sNNNN.tsv}, with four lowercase
hexadecimal digits for the segment identifier; the matching external seal
file contains 64 lowercase hexadecimal digits followed
by LF. The manifest \texttt{index.tsv} begins with
\texttt{A348425-MANIFEST-1}, a TAB, and \texttt{cfg\_sha=} followed by
the configuration digest. Its second line has the TAB-separated labels
\texttt{id}, \texttt{low}, \texttt{high}, \texttt{state\_sha},
\texttt{elapsed\_ns}, \texttt{hits}, in that order.
The manifest column labelled \texttt{state\_sha} stores the digest of
the \emph{entire} segment file, equal to its external seal; it is distinct
from the in-segment field of the same name, which covers only the preceding
segment bytes. Thus the ordered manifest binds the configuration and the
full segment files, including their trailing digest lines.

The aggregate \texttt{hits.tsv} begins with \texttt{A348425-HITS-1},
a TAB, and \texttt{cfg\_sha=} followed by the configuration digest.
Its second line has the TAB-separated column labels \texttt{r},
\texttt{n}, \texttt{x}, \texttt{y}, \texttt{R}, corresponding to the
five coordinates above. Every record line ends in LF.
A completed segment must have $\texttt{next}=\texttt{high}+1$; the next
segment starts at that cursor. The aggregate hit list is regenerated
from the segments and checked for equality, rather than admitted as a
separate authority. These conditions expose missing, reordered or
inconsistent records. They establish integrity and recorded completion;
the exhaustive algorithm remains the basis for negative coverage.

The two saved internal elapsed counters correspond to approximately
1539 and 1387 seconds, respectively; their full precision is retained
in the execution records.
Their sum has the timing limitations stated in Section~\ref{sec:provenance}.
The finite-analysis companion requires SymPy 1.12, deliberately rejecting
other versions to preserve its tested environment. It can be run from an
isolated Python environment with that version installed; it checks positive
rows and class certificates, not exhaustive negative coverage.

\clearpage
\begin{thebibliography}{99}
\sloppy\raggedright
\bibitem{barbeau}
E. J. Barbeau, Remarks on an arithmetic derivative,
\textit{Canadian Mathematical Bulletin} \textbf{4}(2) (1961), 117--122.
\url{https://doi.org/10.4153/CMB-1961-013-0}.

\bibitem{ufn}
V. Ufnarovski and B. \AA hlander, How to differentiate a number,
\textit{Journal of Integer Sequences} \textbf{6} (2003), Article 03.3.4.
\url{https://cs.uwaterloo.ca/journals/JIS/VOL6/Ufnarovski/ufnarovski.html}.

\bibitem{kovic}
J. Kovi\v{c}, The arithmetic derivative and antiderivative,
\textit{Journal of Integer Sequences} \textbf{15} (2012), Article 12.3.8.
\url{https://cs.uwaterloo.ca/journals/JIS/VOL15/Kovic/kovic4.html}.

\bibitem{survey}
T. Tossavainen, P. Haukkanen, J. K. Merikoski, and M. Mattila,
We can differentiate numbers, too,
\textit{The College Mathematics Journal} \textbf{55}(2) (2024), 100--108;
published online November 3, 2023. DOI: 10.1080/07468342.2023.2268494.


\bibitem{book}
G. Balzarotti and P. P. Lava, \textit{La derivata aritmetica:
Alla scoperta di un nuovo approccio alla teoria dei numeri},
Hoepli, Milan, 2013. ISBN 9788820358648.

\bibitem{oeis}
M. A. Burtea, A348425: Squares whose second arithmetic derivative is a
square, \textit{The On-Line Encyclopedia of Integer Sequences},
October 18, 2021; revision 21, September 8, 2022;
snapshot acquired and rechecked October 22, 2026.
Programs by A. Eldar (October 18, 2021) and M. Marcus (October 30, 2021).
\url{https://oeis.org/A348425}.

\bibitem{squares}
N. J. A. Sloane, A000290: The squares: a(n) = n squared,
\textit{The On-Line Encyclopedia of Integer Sequences}.
Snapshot and linked data acquired October 22, 2026.
\url{https://oeis.org/A000290}.

\bibitem{derivative}
N. J. A. Sloane and R. K. Guy, A003415: Arithmetic derivative,
\textit{The On-Line Encyclopedia of Integer Sequences}.
Snapshot and linked data acquired October 22, 2026.
\url{https://oeis.org/A003415}.

\bibitem{second}
R. Zumkeller, A068346: Second arithmetic derivative,
\textit{The On-Line Encyclopedia of Integer Sequences}.
February 28, 2002; table by A. P. Heinz, with initial terms from T. D. Noe.
Snapshot and linked data acquired October 22, 2026.
\url{https://oeis.org/A068346}.

\bibitem{near2}
P. De Geest, A028871: Primes of the form k squared minus 2,
\textit{The On-Line Encyclopedia of Integer Sequences}.
table by N. Johnston.
Snapshot and linked data acquired October 22, 2026.
\url{https://oeis.org/A028871}.

\bibitem{near3}
P. De Geest, A028873: Numbers k such that k squared minus 3 is prime,
\textit{The On-Line Encyclopedia of Integer Sequences}.
table by N. Johnston.
Snapshot and linked data acquired October 22, 2026.
\url{https://oeis.org/A028873}.

\bibitem{near5}
V. Librandi, A201787: Primes of the form 5n squared minus 6,
\textit{The On-Line Encyclopedia of Integer Sequences}.
December 5, 2011; table by V. Librandi.
Snapshot and linked data acquired October 22, 2026.
\url{https://oeis.org/A201787}.

\bibitem{evenpower}
Labos Elemer, A056798: Prime powers with even nonnegative exponents,
\textit{The On-Line Encyclopedia of Integer Sequences}.
August 28, 2000; table by T. D. Noe.
Snapshot and linked data acquired October 22, 2026.
\url{https://oeis.org/A056798}.

\bibitem{powersecond}
R. Zumkeller, A192084: Second arithmetic derivative of squares of prime powers,
\textit{The On-Line Encyclopedia of Integer Sequences}.
June 26, 2011; table by R. Zumkeller.
Snapshot and linked data acquired October 22, 2026.
\url{https://oeis.org/A192084}.

\bibitem{firstsquare}
P. P. Lava, A266890: Squares whose arithmetic derivative is a square,
\textit{The On-Line Encyclopedia of Integer Sequences}.
April 8, 2016; extension and infinitude comment by G. Resta, May 28, 2016.
Snapshot and linked data acquired October 22, 2026.
\url{https://oeis.org/A266890}.

\bibitem{prime2}
V. Shevelev, A192192: Numbers whose second arithmetic derivative is prime,
\textit{The On-Line Encyclopedia of Integer Sequences}.
June 25, 2011; extension by O. G\'erard, July 4, 2011; revised definition by A. Karttunen, October 19, 2019; table by A. Karttunen with initial terms from T. D. Noe.
Snapshot and linked data acquired October 22, 2026.
\url{https://oeis.org/A192192}.

\bibitem{sfree2}
A. Karttunen, A328244: Numbers whose second arithmetic derivative is squarefree,
\textit{The On-Line Encyclopedia of Integer Sequences}.
October 11, 2019.
Snapshot and linked data acquired October 22, 2026.
\url{https://oeis.org/A328244}.

\bibitem{pollard}
J. M. Pollard, A Monte Carlo method for factorization,
\textit{BIT} \textbf{15}(3) (1975), 331--334.
\url{https://doi.org/10.1007/BF01933667}.

\bibitem{miller}
G. L. Miller, Riemann's hypothesis and tests for primality,
\textit{Journal of Computer and System Sciences} \textbf{13}(3)
(1976), 300--317.
\url{https://www.cs.cmu.edu/~glmiller/Publications/b2hd-Mi76.html}.

\bibitem{rabin}
M. O. Rabin, Probabilistic algorithm for testing primality,
\textit{Journal of Number Theory} \textbf{12}(1) (1980), 128--138.
\url{https://www.sciencedirect.com/science/article/pii/0022314X80900840}.

\bibitem{feitsma}
J. Feitsma, with W. Galway (editor), \textit{Tables of pseudoprimes
and related data}, all base-2 Fermat pseudoprimes below $2^{64}$,
last modified April 25, 2013; accessed October 23, 2026.
\url{https://www.cecm.sfu.ca/Pseudoprimes/}.

\bibitem{mr}
W. Izykowski, \textit{Deterministic variants of the Miller--Rabin
primality test}, record of the seven-base set of J. Sinclair (2011),
accessed October 22, 2026.
\url{https://miller-rabin.appspot.com/}.

\bibitem{sympy}
SymPy development team, \textit{SymPy}, version 1.12 (2023),
source release; factorization and primality routines used by the Python
validators. \url{https://github.com/sympy/sympy/releases/tag/sympy-1.12}.

\bibitem{project}
C. Corti, \textit{Squares with Square Second Arithmetic Derivative:
Infinite Families and an Exact Extension of OEIS A348425},
GitHub repository and Zenodo archive, 2026.
Prospective project locators; no public release is asserted in this draft:
\url{https://github.com/carcorti/A348425};
\url{https://doi.org/10.5281/zenodo.xxxxxxxx}.
\end{thebibliography}
\end{document}
